\chapter*{ABSTRAK}
\addcontentsline{toc}{chapter}{ABSTRAK}

\begin{center}

    \textbf{Pengembangan Arsitektur Platform DataOps dan MLOps Terintegrasi pada Kubernetes dengan Pemanfaatan \textit{Open Source Tools}} \\
    \vspace{0.4cm}
    oleh \\
    \vspace{0.4cm}
    Bryan P. Hutagalung\\
    18222130\\
\end{center}


\begin{spacing}{1.0}
Pengembangan model \textit{machine learning} ke tahap produksi masih menghadapi fragmentasi \textit{tool} dan pemisahan kerja antara tim data dan tim model. Praktik DataOps memperkuat kualitas dan kecepatan alur data, sementara MLOps memperkuat keandalan \textit{model lifecycle}, namun keduanya sering tumbuh terpisah sehingga \textit{lineage}, deteksi \textit{drift}, dan \textit{feature serving} antara mode \textit{batch} dan \textit{streaming} tidak konsisten lintas siklus. Berlangganan layanan terkelola memangkas konfigurasi tetapi biaya berjalannya tinggi. Penelitian ini mengembangkan arsitektur platform yang mengintegrasikan DataOps dan MLOps pada Kubernetes dengan komponen \textit{open source}, sehingga satu \textit{control plane} menyatukan sembilan \textit{layer} arsitektur, yaitu \textit{orchestration and infrastructure}, \textit{data ingestion and processing}, \textit{data storage}, \textit{feature service}, \textit{model lifecycle}, \textit{data governance}, \textit{observability}, \textit{GitOps and continuous delivery}, serta \textit{platform security}. Metode \textit{Design Science Research Methodology} (DSRM) mencakup identifikasi masalah, penyusunan tujuan artefak, perancangan, implementasi berbasis GitOps, serta uji coba dengan \textit{use case} pasar aset kripto sebagai kasus verifikasi karena beroperasi tanpa henti. Arsitektur menempatkan Kubernetes sebagai bidang orkestrasi, Apache Kafka dan Apache Flink untuk \textit{streaming}, Apache Spark untuk \textit{batch processing}, Feast sebagai \textit{feature store} ganda dengan ClickHouse \textit{offline}, Valkey \textit{online}, dan Qdrant sebagai indeks vektor, MLflow dan Kubeflow untuk \textit{model lifecycle}, KServe untuk \textit{model serving}, serta DataHub untuk katalog dan \textit{lineage}. Deteksi \textit{drift} berbasis \textit{Population Stability Index} dan Kolmogorov-Smirnov memicu \textit{retraining} otomatis melalui Argo CronWorkflow, sementara \textit{point-in-time correctness} dijaga Feast pada fase \textit{training} dan \textit{online serving}. Argo CD menyatukan GitOps untuk reproduksibilitas \textit{deployment} dari satu sumber Git, sedangkan Prometheus, Grafana, OpenTelemetry, Loki, serta Grafana Tempo melengkapi observabilitas metrik, log, dan \textit{trace}. Hasil evaluasi menunjukkan keempat tujuan platform terpenuhi secara fungsional dan struktural pada lingkungan \textit{node} tunggal, dengan tujuan keempat pada tingkat struktural, mencakup rekonstruksi seluruh komponen dari satu sumber Git, tata kelola otomatis sepanjang \textit{pipeline}, pemicuan \textit{retraining} oleh deteksi \textit{drift}, dan layanan fitur \textit{dual-store} dengan kontrak konsisten antara \textit{training} dan \textit{inference}. Perbandingan biaya tiga skenario \textit{hosting} menunjukkan arsitektur \textit{open source} menekan biaya berjalan jangka panjang dibanding layanan terkelola. Karakterisasi beban kuantitatif pada klaster multi-\textit{node} menjadi arah pengembangan lanjutan sesuai batasan penelitian.

Kata kunci: DataOps, MLOps, Kubernetes, Feature Store, Drift Detection, GitOps, Open Source.
\end{spacing}
